\section{The Prym condition for the symmetric reduction}
\label{sec:geometry-prym-proof}

Under \cref{eq:geometry-fold}, define the constant matrix
\begin{equation}\label{eq:prym-transpose}
 \mathsf S=\begin{pmatrix}
 0&\Id_s&0\\ \Id_s&0&0\\ 0&0&-1
 \end{pmatrix}.
\end{equation}
The matrices in \cref{eq:canonical-X,eq:canonical-lax} obey
\begin{equation}\label{eq:prym-pencil-duality}
 \mathsf S^2=\Id_{2s+1},\qquad
 \mathsf X=\mathsf S\mathsf X^{\trans}\mathsf S,\qquad
 \mathsf S\mathsf D\mathsf S=-\mathsf D,\qquad
 \mathsf T(\nu)=-\mathsf S\mathsf T(-\nu)^{\trans}\mathsf S.
\end{equation}
Consequently
$\mathsf L(\nu)=\mathsf S\mathsf L(-\nu)^{\trans}\mathsf S$,
which proves \cref{eq:geometry-involution}. It also identifies a
right eigenvector at $\iota P$ with a left eigenvector at $P$.

Normalize the last component of the right eigenvector to one, and
let $\mathcal A(t)$ be its finite pole divisor on a generic smooth
spectral chart. Consider the last diagonal cofactor
\begin{equation}\label{eq:prym-cofactor}
 A(\lambda,\nu)=
 \det\bigl(\lambda\Id_{2s}-\mathsf L_{\mathrm{upper}}(\nu)\bigr),
\end{equation}
where $\mathsf L_{\mathrm{upper}}$ is the first $2s$ rows and
columns of the pencil. This cofactor is even in $\nu$.
At a smooth finite point of $C$, the adjugate of
$\lambda\Id_{2s+1}-\mathsf L$ has rank one. In local nonvanishing
eigenvector frames it is a nonzero scalar times a right vector and
a left vector. The zeros of its last diagonal entry thus split
between zeros of their last components. The transpose relation
identifies these two zero divisors by $\iota$, giving
\begin{equation}\label{eq:prym-divisor}
 (A)_0=\mathcal A(t)+\iota\mathcal A(t).
\end{equation}
This is stated on the open chart where these zeros are simple,
disjoint and finite.

Generically, the poles of $A$ have order $s-1$ on each of the $2s$
outer infinity branches and order $2s$ on the central branch.
These orders follow by substituting \cref{eq:geometry-branches}
into \cref{eq:prym-cofactor}. Thus its pole divisor is the fixed divisor
\[
 \mathcal A_\infty
 =(s-1)\left(\sum_{\ell=1}^{s}P_\ell^L
                  +\sum_{m=1}^{s}P_m^M\right)+2s P_1^N,
 \qquad \deg\mathcal A_\infty=2s^2.
\]
It follows that $\deg\mathcal A(t)=s^2$ and
$\mathcal A(t)+\iota\mathcal A(t)\sim\mathcal A_\infty$.
For a fixed reference divisor $\mathcal A(0)$, the degree-zero class
\[
 \xi(t)=[\mathcal A(t)-\mathcal A(0)]\in\operatorname{Jac}(C)
\]
therefore satisfies $\xi(t)+\iota(\xi(t))=0$.
Since $\xi(0)=0$, connected evolution puts $\xi(t)$ in the identity
component, namely the Prym. The standard isospectral
eigenvector-divisor linearization
\cite[\S3]{AdlerVanMoerbeke2018} consequently restricts the motion
to a translate of this Prym. This proves containment of the reduced
divisor motion; a global reconstruction of the entire reduced
isospectral manifold would require its normalization data as well.

For completeness, there are $s$ constant diagonal mode gauges.
Writing $\boldsymbol v=(v_{i1})_{i=1}^{s}$, their action is
\[
 v_{i1}\mapsto e^{\epsilon_i}v_{i1},\qquad
 u_{ij}\mapsto e^{-\epsilon_i-\epsilon_j}u_{ij},
 \qquad 1\leq i,j\leq s,
\]
with inverse factors on the partners and the same constraints on
$\boldsymbol w$. They preserve the normalized eigenvector's pole
divisor. The number of independent pairs after folding is
$s+s(s+1)/2$. Indeed, the pulled-back symplectic form has factor two
on each vector-mode pair and on each off-diagonal symmetric
matrix-mode pair. Multiplying both members of each such pair by
$\sqrt2$, and leaving the diagonal pairs unchanged, gives canonical
coordinates. This accounts for the gauge and Prym dimensions in
\cref{sec:geometry-prym}.
